% ======================================================================
%  Anomaly detection in an industrial robot cell
%  A three-stage evaluation of Vision--Language Models against an
%  off-the-shelf object detector
% ======================================================================
\documentclass[11pt,a4paper]{article}
\usepackage[margin=2.3cm]{geometry}
\usepackage{booktabs}
\usepackage{array}
\usepackage{amssymb}
\usepackage{ragged2e}
\usepackage{tabularx}
\usepackage{placeins}
\usepackage[hidelinks]{hyperref}
\setlength{\parskip}{0.6em}
\setlength{\parindent}{0pt}
\renewcommand{\arraystretch}{1.2}
\newcolumntype{L}[1]{>{\RaggedRight\arraybackslash}p{#1}}

\title{Anomaly Detection in an Industrial Robot Cell:\\[4pt]
       A Three-Stage Evaluation of Vision--Language Models\\
       against an Off-the-Shelf Object Detector}
\author{}
\date{}

\begin{document}
\maketitle

\begin{abstract}
\noindent
This report evaluates whether a vision--language model (VLM) can serve as
the anomaly-detection component of a monitoring system for an industrial
robot cell, and how it compares with a small, off-the-shelf object
detector used for the same purpose. Two categories of anomaly are
considered: \emph{common} anomalies, in which a person or a foreign object
enters the work area, and \emph{operational} anomalies, in which the robot
itself fails to complete its task. Three detectors are compared --- GPT-4o,
Qwen3-VL~8B and Llama~3.2~Vision~11B --- against two rule-based variants of
YOLO11n. Each detector is run through a two-stage procedure (a coarse pass
over the whole clip followed by a dense re-check around the suspected
event) and each of its positive verdicts is then adjudicated by a separate
language model that judges whether the anomaly the detector described is
the one that actually occurred. Every clip is analysed ten times to expose
run-to-run instability, and a further ten-minute recording of normal
operation is used to probe false alarms. The main findings are that all
strong models detect human intrusions reliably but differ sharply on
foreign objects and on operational faults observed without a machine-state
log; that the dense re-check, as currently configured, removes more
correct detections than it adds; and that when a model does raise a false
alarm it is almost always because it has hallucinated a human hand
manipulating the workpieces.
\end{abstract}

\tableofcontents
\newpage

% ======================================================================
\section{Introduction}
% ======================================================================

An automated safety monitor for a robot cell must answer one question
continuously: is the current situation normal, or is something wrong? Two
kinds of ``wrong'' matter in this setting. The first is an
\textbf{external} or \emph{common} anomaly: a human hand, arm or whole
person has entered the robot's work envelope, or an object that has no
business being there --- a tool left behind, a bottle, a ball --- is
sitting in the workspace. These are detectable from the camera alone. The
second is an \textbf{internal} or \emph{operational} anomaly: the robot is
physically intact and unobstructed, but its process has failed --- it
closed its gripper on nothing, it dropped the part, it skipped a pick.
These are only partly visible; in a real installation they are usually
confirmed from the controller's own state log.

The conventional tool for the first category is an object detector such as
YOLO. It is fast, deterministic and cheap, but it is a closed-vocabulary
classifier: it can only report the categories it was trained on, and it
has no notion of ``this object does not belong here''. A vision--language
model is the opposite trade-off. It reasons in open vocabulary, it can be
told the operational context in words, and it can read a state log
alongside the images; but it is slow, non-deterministic, and it can
hallucinate.

This report presents a controlled comparison of the two approaches across
three tests: detection of common anomalies (Section~\ref{sec:test1}),
detection of operational anomalies with and without the state log
(Section~\ref{sec:test2}), and false-alarm behaviour on a long recording
of normal operation (Section~\ref{sec:test3}). Section~\ref{sec:method}
describes the detection pipeline, the semantic-adjudication step, the
detectors and the test material. Section~\ref{sec:discussion} draws the
comparison together.

% ======================================================================
\section{Evaluation methodology}
\label{sec:method}
% ======================================================================

\subsection{The two-stage detection pipeline}

Every detector produces its verdict in two stages.

\textbf{Stage~1 --- coarse pass.} The clip is sampled at nine points
spaced evenly across its full duration (frame $\lfloor i(N-1)/8\rfloor$
for $i = 0,\dots,8$, where $N$ is the clip's frame count). Each of the
nine frames is cropped to the work surface using a colour heuristic,
scaled so that its longer side is at most 560 pixels, and JPEG-encoded.
For operational-anomaly clips the machine-state log, when one is
available, is appended to the prompt as plain text. The detector receives
all nine frames together and must return a short description of each frame
followed by a single verdict line, either
\texttt{ANOMALY:}~\emph{justification} or \texttt{NORMAL:}~\emph{justification}.
From the justification the pipeline extracts the \emph{anchor}: the frame
number, or numbers, the detector says the anomaly is in. If the detector
produces no parseable verdict line at all --- which happens when a model
refuses or drifts into an unstructured monologue --- the verdict is
recorded as \texttt{UNKNOWN}.

The nine-frame coarse pass is deliberately sparse. A brief event --- a
hand that enters and withdraws within a second, a ball that crosses the
frame --- can fall between two of the nine samples, or be caught in only
one, motion-blurred frame. Stage~1 is therefore expected to be sensitive
but imprecise about \emph{when} the anomaly occurs.

\textbf{Stage~2 --- dense re-check.} Stage~2 runs only when Stage~1
returned \texttt{ANOMALY} and an anchor frame could be identified; a
Stage~1 verdict of \texttt{NORMAL} or \texttt{UNKNOWN} is never revisited.
Around each anchor the pipeline builds a dense window of frames --- the
anchor plus and minus multiples of four source frames, up to ten steps in
each direction, then de-duplicated and capped at twelve frames for the
language models. These frames are prepared exactly as in Stage~1, the log
is appended again, and the \emph{same} detector is called afresh with the
\emph{same} prompt. Crucially, the model is not shown its own Stage~1
answer; Stage~2 is an independent second opinion on a tighter, denser view
of the same moment.

The Stage~2 verdict is taken as the \textbf{final} verdict. It can
therefore \emph{overturn} Stage~1, turning an \texttt{ANOMALY} into a
\texttt{NORMAL}. The reverse cannot happen, because Stage~2 never runs on
a Stage~1 \texttt{NORMAL}. If Stage~2 itself returns \texttt{UNKNOWN} or
errors, the Stage~1 \texttt{ANOMALY} is retained. Stage~2 also produces
the localised time window of the anomaly --- the first and last timestamps
in the dense window at which the detector still reports the event.

The design intent of Stage~2 is twofold: to confirm a suspected event
against evidence the sparse pass could not have seen, and to pin down its
timing. On a test set in which every clip genuinely contains an anomaly,
an overturning Stage~2 can only lower recall; its benefit is expected to
appear instead on normal footage, where it has the opportunity to retract
a false alarm.

\subsection{Semantic adjudication (Stage~3)}

A verdict of \texttt{ANOMALY} can be correct as a label while being wrong
in substance: the model may have flagged the right clip for the wrong
reason, or described a plausible-sounding anomaly that is not the one
present. To separate these cases we add a third stage that does not affect
any verdict and is applied only after all runs are complete.

For every run whose \emph{final} verdict is \texttt{ANOMALY}, an
independent language model (GPT-4o at temperature zero) is given two short
texts: a one-sentence reference description of the true anomaly for that
clip, and the detector's own final justification. It returns exactly one
of three labels. \textbf{MATCH} means the justification identifies the
same anomaly type and the same agent as the reference --- a human
hand or arm, a specific foreign object such as a ball, bottle, plant or
wallet, or the same operational fault such as the gripper closing on
nothing. Different wording or a missing minor detail does not prevent a
MATCH. \textbf{PARTIAL} means the core is right but the account is
incomplete or vague, or it is right but padded with a clearly false extra
claim. \textbf{MISMATCH} means the justification describes a different
anomaly, one that was hallucinated, or --- on the normal video --- any
anomaly at all. Runs whose final verdict is \texttt{NORMAL} or
\texttt{UNKNOWN} are treated as ``no detection'' and are not sent to the
judge. Identical justifications are judged once and the result reused; the
618 positive runs reduce to 362 distinct texts.

Stage~3 lets us report not only \emph{how often} a detector fires on a
real anomaly, but \emph{how often its explanation is trustworthy}.

\subsection{Detectors under evaluation}

\textbf{GPT-4o} is used through the OpenAI API at temperature zero. It
accepts all nine (or twelve) frames in a single request and returns the
description and verdict directly.

\textbf{Qwen3-VL~8B} is run locally through Ollama, also receiving all
frames in one request. It is a ``thinking'' model: left to itself it
prefixes its answer with an extended internal monologue. It is therefore
forced into non-thinking mode and given a generous output budget, but even
so it occasionally produces an essay with no verdict line, which the
pipeline records as \texttt{UNKNOWN}.

\textbf{Llama~3.2~Vision~11B} is run locally through Ollama in a
four-bit quantisation. Unlike the other two it accepts only one image per
request, so it cannot be given the frames together. Instead each frame is
captioned individually and the nine captions are then passed as text to a
final synthesis call that produces the verdict. This indirection is the
model's main handicap: fine visual detail that is not mentioned in a
caption is lost before the decision is made.

\textbf{YOLO11n} and \textbf{YOLO11n-pose} are the two rule-based
baselines. The detection variant runs the standard 80-class model on each
frame and declares an anomaly if the class \texttt{person} is found at
confidence at least $0.35$ in any frame. The pose variant declares an
anomaly if a wrist keypoint is visible at the same confidence in any
frame, which is a finer proxy for a hand reaching into the cell. Both are
deterministic; their ten repeats per clip are identical by construction
and are reported at $n=10$ only so that the tables line up with the
language models. In this system YOLO is normally used only to crop the
region of interest; here it is pressed into service as an actual anomaly
detector through the rules above.

\subsection{Test material}

The common-anomaly set consists of ten short clips (three to eighteen
seconds, recorded at 60\,fps). Seven contain a human hand or arm entering
the workspace; three contain a foreign object only --- a thrown tennis
ball, a potted plant, and a plastic bottle that the robot arm then
collides with; one contains both a hand and the removal of a workpiece.

The operational-anomaly set is built to isolate the contribution of the
state log. Five clips (an1--an5) show a process fault --- the robot picks
a piece that is not there, picks nothing, or fails to grasp --- and are
accompanied by the controller's state log, which records the gripper
error explicitly. The other five clips (an6--an10) are \emph{the same five
recordings, byte for byte, with the log removed}. The difference between
the two halves is therefore purely the presence of the log, and it
measures how much of the operational-anomaly detection depends on that
text rather than on the images.

Neither anomaly set contains any clip of normal operation. All reported
figures on those sets are therefore \emph{recall} --- the fraction of
anomalies detected --- and cannot speak to precision. To probe the
false-positive side we use a separate ten-minute recording
(\texttt{teste\_falso\_alarm.mp4}, 14\,412 frames at 24\,fps) of
uninterrupted normal operation, in which every \texttt{ANOMALY} verdict is
by definition an error.

\subsection{Protocol and run counts}

Each of the three language models was run ten times on each of the twenty
anomaly clips, giving 600 language-model runs; the two YOLO variants were
run (deterministically) on the ten external clips, giving 200 rows. On the
false-alarm video each language model was run six times, and each of those
runs drew a fresh random set of frames: a random count between nine and
eighteen, at random timestamps across the whole recording, with a distinct
set per model and run, so that the same recording is never analysed twice
through the same frames. This yields 18 false-alarm runs. Stage~3 produced
618 judgements. Repeating every anomaly clip ten times is what makes the
run-to-run instability of the language models visible; a model that
detects a clip in six runs out of ten is reported as ``6/10'', not simply
as ``detected''.

% ======================================================================
\section{Test 1 --- Common (external) anomalies}
\label{sec:test1}
% ======================================================================

\subsection{Anomalies Test Results}

The ten common-anomaly clips were analysed by all five detectors, ten
times each. For the language models we report, per clip, how many of the
ten runs raised an anomaly at Stage~1 and how many still did after the
Stage~2 re-check. For YOLO the two numbers are always equal. Across the
whole set we then report the Stage~1 recall, the final recall, the
fraction of clips in which the correct agent was named (the coarse
right-reason check), and the Stage~3 verdict on the justifications.

\subsection{Stage 1: Coarse Detection}

The ten common-anomaly clips were analysed by all five detectors, ten
times each. For the language models we report, per clip, how many of the
ten runs raised an anomaly at Stage~1. The results can be seen on
Table~\ref{tab:test1-stage1}.

\begin{table}[h]
\centering
\small
\setlength{\tabcolsep}{5pt}
\begin{tabular}{@{}l ccc cc@{}}
\toprule
Clip & GPT-4o & Qwen3-VL & Llama-V & YOLO-coco & YOLO-pose \\
\midrule
an1  & 10 & 10 & 10 & 0  & 10 \\
an2  & 10 & 10 & 10 & 10 & 10 \\
an3  & 10 & 10 & 10 & 10 & 10 \\
an4  & 10 & 10 & 10 & 10 & 10 \\
an5  & 10 & 10 & 0  & 0  & 0 \\
an6  & 10 & 10 & 10 & 10 & 10 \\
an7  & 10 & 10 & 10 & 10 & 10 \\
an8  & 10 & 10 & 10 & 10 & 10 \\
an9  & 10 & 0  & 0  & 0  & 10 \\
an10 & 10 & 10 & 10 & 10 & 10 \\
\midrule
\textbf{recall} & \textbf{100\,\%} & \textbf{90\,\%} & \textbf{80\,\%}
                & \textbf{70\,\%} & \textbf{90\,\%} \\
\midrule
mean time (s)   & 8.8 & 36.2 & 71.1 & $\approx$2 & $\approx$2 \\
median time (s) & 6.3 & 35.7 & 80.0 & $\approx$1.5 & $\approx$1.5 \\
\bottomrule
\end{tabular}
\caption{Test~1 (common / external anomalies), Stage~1 coarse pass:
number of \texttt{ANOMALY} verdicts out of ten runs per clip, for each
detector, with per-run wall-clock time (Stage~1 and Stage~2 combined).}
\label{tab:test1-stage1}
\end{table}
\FloatBarrier

At Stage 1, GPT-4o raised an anomaly in every one of its hundred runs
(ten out of ten on all ten clips). Qwen3-VL matched this on nine of the
ten clips and missed the plastic-bottle clip (an9) entirely, giving 90 of
100. YOLO's pose variant also reached 90 of 100, failing only on the
tennis-ball clip (an5), where there is no person and no visible wrist.
Llama reached 80 of 100, since like Qwen it never fired on the bottle and
it additionally never fired on the ball. YOLO's plain detection variant
reached 70 of 100, missing the ball, the bottle, and, because the hand
never forms a full human silhouette, the brief hand in an1.

The pattern at this stage is that human intrusions are easy to detect. On
the seven clips that contain a hand or arm (an1 to an4, an6 to an8, an10),
every detector fired in all ten runs, with the single exception of YOLO's
detection variant on the one-frame hand in an1. The difficulty is
concentrated in the three object-only clips. The thrown tennis ball (an5)
was detected by GPT-4o and Qwen in all ten runs, but by neither YOLO
variant and never by Llama. The static plastic bottle (an9) was detected
only by GPT-4o, in all ten Stage-1 runs, and spuriously by YOLO-pose,
which fired a wrist keypoint on a scene that contains only the bottle and
the robot arm. The potted plant (an6) was detected by all three language
models in every run, and also by both YOLO variants, but YOLO fired for
the wrong reason, reading the cluttered scene as containing a person
rather than recognizing the plant as a foreign object.

\subsection{Stage 2: the dense re-check}

\begin{table}[h]
\centering
\small
\setlength{\tabcolsep}{5pt}
\begin{tabular}{@{}l ccc cc@{}}
\toprule
Clip & GPT-4o & Qwen3-VL & Llama-V & YOLO-coco & YOLO-pose \\
\midrule
an1  & 10 & 10 & \textbf{5} & 0  & 10 \\
an2  & 10 & 10 & 10 & 10 & 10 \\
an3  & 10 & 10 & 10 & 10 & 10 \\
an4  & 10 & 10 & 10 & 10 & 10 \\
an5  & \textbf{6} & 10 & 0 & 0 & 0 \\
an6  & 10 & 10 & 10 & 10 & 10 \\
an7  & 10 & 10 & \textbf{8} & 10 & 10 \\
an8  & 10 & 10 & 10 & 10 & 10 \\
an9  & \textbf{3} & 0 & 0 & 0 & 10 \\
an10 & 10 & \textbf{0} & 10 & 10 & 10 \\
\midrule
\textbf{final recall} & \textbf{89\,\%} & \textbf{80\,\%} & \textbf{73\,\%}
                      & \textbf{70\,\%} & \textbf{90\,\%} \\
\midrule
Stage-1 recall (ref.) & 100\,\% & 90\,\% & 80\,\% & 70\,\% & 90\,\% \\
runs overturned & 11 & 10 & 7 & 0 & 0 \\
\bottomrule
\end{tabular}
\caption{Test~1, Stage 2 results.}
\label{tab:test1-stage2}
\end{table}
\FloatBarrier

On this test set every clip is an anomaly, so any change Stage~2 makes to
the verdict is a change for the worse. Stage~2 changed the verdict in 23
of the 300 language-model runs, and every one of those 23 changes was an
\texttt{ANOMALY} turned into a \texttt{NORMAL}, resulting in twenty-three
lost detections.

The losses were not spread evenly. They clustered on a handful of
clip--model combinations, each with an identifiable mechanism. GPT-4o lost
seven of its ten runs on the bottle clip (an9) because Stage~2, shown the
bottle clearly in every frame of the dense window, re-reasoned about it
and concluded that the bottle does not interfere with the operation and
does not represent a safety anomaly, a defensible-sounding judgment that
nonetheless contradicts the task definition, under which any object that
does not belong is an anomaly. It lost a further four of ten on the ball
clip (an5) for a different reason. The anchor extracted from Stage~1
pointed at the middle of the clip, so the dense window covered roughly
0.9 to 2.1 seconds, which is after the ball has crossed the frame (0.4 to
0.8 seconds). Stage~2 then truthfully reported that it saw only the robot
arms and the yellow blocks, and returned \texttt{NORMAL}. This second case
is not a reasoning failure at all, therefore it is a sampling failure in
how the dense window is placed.

Qwen3-VL lost all ten of its runs on the mixed clip an10 to a generation
failure. On every one of those runs its Stage~2 output was the unfilled
prompt template itself, the literal string
``\texttt{NORMAL: <factual and elaborate explanation\ldots>}'', which the
parser dutifully read as a \texttt{NORMAL} verdict. This single failure is
the entire difference between Qwen's 90\,\% Stage-1 recall and its 80\,\%
final recall on this test. On a further ten runs, all on clip an2, Qwen's
Stage~2 returned an unstructured monologue with no verdict line. Here the
guard did its job, the Stage~1 \texttt{ANOMALY} was kept, and the only
cost was ten wasted Stage-2 calls.

Llama lost five of ten runs on the one-frame hand in an1, since it was not
capable of detecting the frame containing the hand, and Stage~2, working
from captions of frames that show only the workbench, described ``two
microscopes on a blue table'' and returned \texttt{NORMAL}. It lost two
more on an7 to a scene-hallucination flip, in which Stage~1 reported ``a
hand in frame $X$'' and Stage~2, re-describing the same scene, asserted
that ``no people are present and welding proceeds normally''.

\subsection{Stage 3: quality of the justifications}

Of the runs that ended as \texttt{ANOMALY}, Stage~3 asked how many
actually described the anomaly that was present. The references for this
comparison can be seen in Table~\ref{tab:test1-stage3-start}.

\begin{table}[htbp]
\centering
\caption{External anomaly test scenarios.}
\label{tab:external_anomaly_scenarios}
\begin{tabularx}{\textwidth}{@{}l l X@{}}
\toprule
\textbf{Test} & \textbf{Anomaly type} & \textbf{Scenario description} \\
\midrule
AN1 & External & A human hand places a white object in the workspace and then withdraws. \\
AN2 & External & A human hand briefly enters the workspace and then withdraws. \\
AN3 & External & A human hand enters the workspace and touches the robot. \\
AN4 & External & A human hand enters the workspace and rearranges the yellow parts. \\
AN5 & External & A tennis ball is thrown into the workspace. \\
AN6 & External & A plant is placed in the workspace. \\
AN7 & External & A human hand places a white object in the workspace. \\
AN8 & External & A human hand throws a black wallet into the workspace while a white object is also present. \\
AN9 & External & A plastic bottle is present in the workspace. \\
AN10 & External & A human hand forcibly removes a part from the robot. \\
\bottomrule
\label{tab:test1-stage3-start}
\end{tabularx}
\end{table}
\FloatBarrier

The most striking result is a positive one: not a single run on this test
set was graded MISMATCH. When a language model fired on a clip that
genuinely contained an anomaly, it never invented a different anomaly. At
worst, it simply under-described the real one.

\begin{table}[h]
\centering
\small
\setlength{\tabcolsep}{7pt}
\begin{tabular}{@{}l ccc@{}}
\toprule
Clip & GPT-4o & Qwen3-VL & Llama-V \\
\midrule
an1  & 1/10  & 8/10  & 0/5  \\
an2  & 8/10  & 0/10  & 1/10 \\
an3  & 10/10 & 10/10 & 0/10 \\
an4  & 10/10 & 10/10 & 2/10 \\
an5  & 6/6   & 10/10 & ---  \\
an6  & 10/10 & 10/10 & 7/10 \\
an7  & 9/10  & 0/10  & 0/8  \\
an8  & 10/10 & 0/10  & 3/10 \\
an9  & 3/3   & ---   & ---  \\
an10 & 0/10  & ---   & 0/10 \\
\midrule
\textbf{MATCH (of fired)} & \textbf{67/89 (75\,\%)} & \textbf{48/80 (60\,\%)}
                          & \textbf{13/73 (18\,\%)} \\
\midrule
PARTIAL   & 22 & 32 & 60 \\
MISMATCH  & 0  & 0  & 0 \\
\bottomrule
\end{tabular}
\caption{Test 1 Stage 3 Results.}
\label{tab:test1-stage3}
\end{table}
\FloatBarrier

Within the positives, however, the models separate widely. GPT-4o's
justifications were graded MATCH on 67 of its 89 positive runs. The
remainder were PARTIAL, concentrated on the composite clip an10, where GPT
consistently reported ``a hand in the work area'' but never mentioned that
the hand was removing a workpiece, so the action that makes an10 distinct.
Qwen3-VL reached MATCH on 48 of 80. It is reliable on hands and on the
plant, but it describes the specific objects generically. On the
white-object clip (an7) and the wallet clip (an8) it fired in all ten runs
yet was graded MATCH in none, because it said only ``an object'' where the
reference names a particular one. Llama reached MATCH on just 13 of its 73
positive runs. It almost always mentions ``a hand'' somewhere, but that
mention is embedded in a wholesale hallucination of the scene, by stating
there were recurring ``microscopes on a blue table''. The judge grades the
account PARTIAL, since the model has noticed that something is wrong
without describing the situation.

The YOLO variants provide no free-text justification, so their detections
were checked manually. Human-intrusion detections were consistent with the
scene, but an6 and an9 produced false person or wrist detections,
classified as \texttt{MISMATCH}. These cases show that detecting an anomaly
does not necessarily mean correctly identifying its cause.

\subsection{Conclusion for Test 1}

The entire results of Test 1 can be seen in Table~\ref{tab:test1}.

\begin{table}[h]
\centering
\begin{tabular}{@{}l c c c c@{}}
\toprule
Detector & Stage-1 recall & Final recall & Right-reason & Stage-3 MATCH$^{\ast}$ \\
\midrule
GPT-4o          & 100/100 (100\,\%) & 89/100 (89\,\%) & 8.9/10 & 67/89 (75\,\%) \\
Qwen3-VL 8B     & 90/100  (90\,\%)  & 80/100 (80\,\%) & 8.0/10 & 48/80 (60\,\%) \\
Llama-V 11B     & 80/100  (80\,\%)  & 73/100 (73\,\%) & 7.0/10 & 13/73 (18\,\%) \\
YOLO11n-pose    & 90/100  (90\,\%)  & 90/100 (90\,\%) & 7.0/10 & --- \\
YOLO11n-coco    & 70/100  (70\,\%)  & 70/100 (70\,\%) & 6.0/10 & --- \\
\bottomrule
\end{tabular}
\caption{Test~1, Final Results.}
\label{tab:test1}
\end{table}
\FloatBarrier

On common anomalies the camera-only task divides cleanly into two regimes.
Human intrusions are effectively solved. Every detector, including both
YOLO variants, catches them in almost every run, and the only lapse is
YOLO's detection variant on a hand too small and too brief to register as
a person. Foreign objects are where the approaches diverge. Only GPT-4o
detected the static bottle, and only GPT-4o and Qwen3-VL detected the
thrown ball. The YOLO rules cannot see objects at all unless a person
happens to be in frame, and Llama's caption pipeline discards them before
the decision. This is the expected consequence of open- versus
closed-vocabulary reasoning, and it is the main argument for a language
model in this role.

The two-stage design did not help here and cost real detections.
Twenty-three correct anomalies were retracted at Stage~2 and none were
gained, because the test set has no normal clips on which a retraction
could be correct. Roughly half of the retractions were genuine
re-judgements, with the model looking again and talking itself out of a
correct call, while the rest were mechanical, a broken model output parsed
as \texttt{NORMAL}, or a dense window placed over the wrong part of the
clip.

According to Test 1, the GPT-4o-based configuration of \textit{Robot
Monitor} achieved a right-reason score of 8.9/10 (89\,\%), outperforming
YOLO11n-pose at 7.0/10 (70\,\%) and YOLO11n-coco at 6.0/10 (60\,\%).
Although YOLO11n-pose achieved a slightly higher final recall (90\,\%
versus 89\,\%), its detections were less consistently associated with the
actual cause of the anomaly.

Furthermore, GPT-4o provided contextual explanations of the detected
events, whereas YOLO outputs were limited to object classes or keypoints.
The GPT-4o configuration therefore offered both stronger identification of
the correct anomaly cause and richer information for interpreting the
event. However, explanation quality was not uniform, since 67 of its 89
positive outputs (75\,\%) received a Stage-3 \texttt{MATCH}.

Detecting an anomaly and describing it correctly are different abilities.
GPT-4o did both on three quarters of its positive runs, while Qwen3-VL
described roughly three fifths of its positives adequately but tends to
name objects only as ``an object''. Llama, although it fired on most
clips, produced a trustworthy account on fewer than one positive run in
five, because its detections are wrapped in a hallucinated scene.

% ======================================================================
\section{Test 2: Operational Internal Anomalies}
\label{sec:test2}
% ======================================================================

\subsection{Anomalies Test Results}

The ten operational-anomaly clips were analysed by the three language
models, ten times each. YOLO was not run here, because the faults have no
visual signature that its rules could key on and it has no way to read the
robot state log. Five of the clips (an1 to an5) were presented with the
state log, and the other five (an6 to an10) are the same five recordings
with the log removed. For each model we report, per clip, how many of the
ten runs raised an anomaly at Stage~1 and how many still did after the
Stage~2 re-check. Across the whole set we then report the Stage~1 recall,
the final recall, the fraction of clips in which the correct cause was
named (the coarse right-reason check), and the Stage~3 verdict on the
justifications.

\subsection{Stage 1: Coarse Detection}

The ten operational-anomaly clips were analysed by the three language
models, ten times each. For each model we report, per clip, how many of
the ten runs raised an anomaly at Stage~1. The results can be seen on
Table~\ref{tab:test2-stage1}.

\begin{table}[h]
\centering
\small
\setlength{\tabcolsep}{7pt}
\begin{tabular}{@{}l ccc@{}}
\toprule
Clip & GPT-4o & Qwen3-VL & Llama-V \\
\midrule
an1  & 10 & 10 & 10 \\
an2  & 10 & 10 & 10 \\
an3  & 10 & 10 & 10 \\
an4  & 10 & 10 & 10 \\
an5  & 10 & 10 & 10 \\
an6  & 7  & 10 & 10 \\
an7  & 10 & 10 & 10 \\
an8  & 10 & 10 & 0 \\
an9  & 10 & 0  & 0 \\
an10 & 9  & 10 & 1 \\
\midrule
\textbf{recall} & \textbf{96\,\%} & \textbf{90\,\%} & \textbf{71\,\%} \\
\midrule
mean time (s)   & 10.4 & 41.1 & 68.2 \\
median time (s) & 6.7  & 44.6 & 81.4 \\
\bottomrule
\end{tabular}
\caption{Test~2 (operational / internal anomalies), Stage~1 coarse pass:
number of \texttt{ANOMALY} verdicts out of ten runs per clip, for each
model, with per-run wall-clock time (Stage~1 and Stage~2 combined). an1 to
an5 include the state log, an6 to an10 do not.}
\label{tab:test2-stage1}
\end{table}
\FloatBarrier

At Stage 1, every model raised an anomaly in all ten runs of the five
clips that include the state log, so on that half of the set GPT-4o,
Qwen3-VL and Llama all reached 50 of 50. This is expected, since the log
contains an explicit gripper error and the prompt tells the model that a
non-null error is always an anomaly. On the five clips without the log the
models separate. GPT-4o reached 46 of 50, missing three runs on an6 and
one on an10, and the an10 miss was an API error rather than a real
\texttt{NORMAL}. Qwen3-VL reached 40 of 50, catching an6, an7, an8 and
an10 in every run but returning \texttt{UNKNOWN} on all ten runs of an9.
Llama reached only 21 of 50, since it never fired on an8 or an9 and fired
only once on an10. Over the whole ten-clip set this gives GPT-4o 96 of
100, Qwen3-VL 90 of 100, and Llama 71 of 100.

The pattern at this stage is that the state log makes the task almost
trivial. When the decisive evidence is handed to the model in text, all
three detect the fault in every run. The difficulty is concentrated in the
five clips where nothing but the robot's own motion is visible. Clip an7
(no log) shows a person visibly reaching in and removing the workpiece,
and every model fired on it in all ten runs, exactly as if it were a
common anomaly, because the human hand is plainly in frame. The genuine
test is the three clips in which only the robot moves, an6 and an8 (the
robot closes on a piece that is not there) and an9 (the robot skips the
pick entirely). GPT-4o and Qwen handled an6 and an8 well but diverged on
an9, where GPT fired in six of ten runs and Qwen in none. Llama fired on
an8 in zero runs and on an10 in one, so without the log it recovers barely
a fifth of the faults.

\subsection{Stage 2: the dense re-check}

\begin{table}[h]
\centering
\small
\setlength{\tabcolsep}{7pt}
\begin{tabular}{@{}l ccc@{}}
\toprule
Clip & GPT-4o & Qwen3-VL & Llama-V \\
\midrule
an1  & 10 & 10 & 10 \\
an2  & 10 & 10 & 10 \\
an3  & 10 & 10 & 10 \\
an4  & 10 & 10 & 10 \\
an5  & 10 & 10 & 10 \\
an6  & \textbf{5} & 10 & \textbf{0} \\
an7  & 10 & 10 & 10 \\
an8  & 10 & 10 & 0 \\
an9  & \textbf{6} & 0 & 0 \\
an10 & 9 & 10 & 1 \\
\midrule
\textbf{final recall} & \textbf{90\,\%} & \textbf{90\,\%} & \textbf{61\,\%} \\
\midrule
Stage-1 recall (ref.) & 96\,\% & 90\,\% & 71\,\% \\
runs overturned & 6 & 0 & 10 \\
\bottomrule
\end{tabular}
\caption{Test~2, Stage 2 results.}
\label{tab:test2-stage2}
\end{table}
\FloatBarrier

On this test set every clip is an anomaly, so any change Stage~2 makes to
the verdict is a change for the worse. Stage~2 changed the verdict in 16
of the 300 language-model runs, and every one of those 16 changes was an
\texttt{ANOMALY} turned into a \texttt{NORMAL}, resulting in sixteen lost
detections. All sixteen fall on the clips without the log, so with the log
Stage~2 never disturbed a correct verdict.

The losses were not spread evenly. They clustered on a handful of
clip--model combinations, each with an identifiable mechanism. Llama lost
all ten of its runs on the phantom-pick clip an6 because its Stage~2
synthesis fabricated a state log, asserting that the state log records no
anomaly and the gripper exerts consistent pressure, for a clip that has no
log at all, and then used that invention to justify \texttt{NORMAL}. This
single failure is the entire difference between Llama's 71\,\% Stage-1
recall and its 61\,\% final recall on this test.

GPT-4o lost four of its ten runs on the no-pick clip an9 and two of its
ten runs on an6 to a visual flip-flop on ambiguous footage. On those runs
Stage~1 reported that the robot grips nothing, which suggests a fault, and
Stage~2, looking again at the same frames, decided that the robot appears
to operate normally and grips the object, and returned \texttt{NORMAL}.
This is not a mechanical failure, therefore it is a genuine second-thought
that talked the model out of a correct call. On a further nine runs, five
on an10 and four on an2, GPT-4o's Stage~2 returned an unstructured
monologue with no verdict line. Here the guard did its job, the Stage~1
\texttt{ANOMALY} was kept, and the only cost was nine wasted Stage-2
calls.

Qwen3-VL produced no internal overturns at all. On the ninety runs where
it raised an anomaly at Stage~1, Stage~2 confirmed the verdict every time.

\subsection{Stage 3: quality of the justifications}

Of the runs that ended as \texttt{ANOMALY}, Stage~3 asked how many
actually described the anomaly that was present. The references for this
comparison can be seen in Table~\ref{tab:test2-stage3-start}.

\begin{table}[htbp]
\centering
\caption{Internal anomaly test scenarios.}
\label{tab:internal_anomaly_scenarios}
\begin{tabularx}{\textwidth}{@{}l l X@{}}
\toprule
\textbf{Test} & \textbf{Anomaly type} & \textbf{Scenario description} \\
\midrule
AN1 & Internal & The robot picks a part that is not present. The state log is available. \\
AN2 & Internal & A human hand forcibly removes a part from the robot. The state log is available. \\
AN3 & Internal & The robot picks a part that is not present. The state log is available. \\
AN4 & Internal & The robot does not pick any part. The state log is available. \\
AN5 & Internal & The robot fails to grasp the part. The state log is available. \\
AN6 & Internal & The robot picks a part that is not present. No state log is provided. \\
AN7 & Internal & A human hand forcibly removes a part from the robot. No state log is provided. \\
AN8 & Internal & The robot picks a part that is not present. No state log is provided. \\
AN9 & Internal & The robot does not pick any part. No state log is provided. \\
AN10 & Internal & The robot fails to grasp the part. No state log is provided. \\
\bottomrule
\label{tab:test2-stage3-start}
\end{tabularx}
\end{table}
\FloatBarrier

The most striking result is a positive one: not a single run on this test
set was graded MISMATCH. When a language model fired on a clip that
genuinely contained an anomaly, it never invented a different anomaly. At
worst, it simply under-described the real one.

\begin{table}[h]
\centering
\small
\setlength{\tabcolsep}{7pt}
\begin{tabular}{@{}l ccc@{}}
\toprule
Clip & GPT-4o & Qwen3-VL & Llama-V \\
\midrule
an1  & 9/10  & 10/10 & 8/10 \\
an2  & 8/10  & 10/10 & 2/10 \\
an3  & 10/10 & 10/10 & 9/10 \\
an4  & 10/10 & 10/10 & 7/10 \\
an5  & 10/10 & 10/10 & 7/10 \\
an6  & 0/5   & 10/10 & ---  \\
an7  & 1/10  & 10/10 & 1/10 \\
an8  & 9/10  & 10/10 & ---  \\
an9  & 5/6   & ---   & ---  \\
an10 & 8/9   & 10/10 & 0/1  \\
\midrule
\textbf{MATCH (of fired)} & \textbf{70/90 (78\,\%)} & \textbf{90/90 (100\,\%)}
                          & \textbf{34/61 (56\,\%)} \\
\midrule
PARTIAL   & 20 & 0 & 27 \\
MISMATCH  & 0  & 0 & 0 \\
\bottomrule
\end{tabular}
\caption{Test 2 Stage 3 Results.}
\label{tab:test2-stage3}
\end{table}
\FloatBarrier

Within the positives, however, the models separate widely. Qwen3-VL's
justifications were graded MATCH on all 90 of its 90 positive runs, with
no PARTIAL and no MISMATCH, including the no-log clips an6, an8 and an10,
on which it consistently stated that the robot completed its motion but
the gripper never closed on a part. GPT-4o reached MATCH on 70 of 90. It
is strong on the clips with the log, but its no-log positives are often
PARTIAL, because on an7 it reports the visible hand rather than the fault,
and because its wording on the phantom-pick clips is frequently vague.
Llama reached MATCH on 34 of 61. With the log it merely echoes the error
text without connecting it to the frames, and without the log it rarely
fires at all.

YOLO was not part of this test, since the operational faults have no
visual signature it could detect and it cannot read the state log. The
comparison here is therefore between the three language models only.

\subsection{Conclusion for Test 2}

The entire results of Test 2 can be seen in Table~\ref{tab:test2}.

\begin{table}[h]
\centering
\begin{tabular}{@{}l c c c c@{}}
\toprule
Detector & Stage-1 recall & Final recall & Right-reason & Stage-3 MATCH$^{\ast}$ \\
\midrule
GPT-4o          & 96/100 (96\,\%) & 90/100 (90\,\%) & 6.8/10 & 70/90 (78\,\%) \\
Qwen3-VL 8B     & 90/100 (90\,\%) & 90/100 (90\,\%) & 7.0/10 & 90/90 (100\,\%) \\
Llama-V 11B     & 71/100 (71\,\%) & 61/100 (61\,\%) & 6.0/10 & 34/61 (56\,\%) \\
\bottomrule
\end{tabular}
\caption{Test~2, Final Results.}
\label{tab:test2}
\end{table}
\FloatBarrier

On operational anomalies the task divides cleanly into two regimes. With
the state log the problem is effectively solved. Every model detects the
fault in all ten runs of all five clips, so the with-log half does not
separate the detectors at all. Without the log the approaches diverge.
GPT-4o and Qwen3-VL both recover about four fifths of the faults from the
images alone, while Llama recovers barely a fifth, since its
caption-then-synthesise pipeline loses the fine cue that the gripper is
open and empty at the end of a motion. The single fault that no model
detects without the log is the one in which the robot skips the pick
entirely, because it produces almost no visual evidence, which is the
point at which the state log stops being a convenience and becomes a
requirement.

The two-stage design did not help here and cost real detections. Sixteen
correct anomalies were retracted at Stage~2 and none were gained, because
the test set has no normal clips on which a retraction could be correct.
Ten of the retractions were Llama on a single clip, where its Stage~2 step
invented a clean log and reasoned from it, while the other six were GPT-4o
flip-flopping on the two hardest no-log clips. Qwen3-VL made no
retractions at all.

According to Test 2, the Qwen3-VL-based configuration of \textit{Robot
Monitor} achieved a right-reason score of 7.0/10 (70\,\%), slightly ahead
of GPT-4o at 6.8/10 (68\,\%) and Llama at 6.0/10 (60\,\%). GPT-4o and
Qwen3-VL reached the same 90\,\% final recall, but Qwen's detections were
more consistently associated with the actual cause of the fault.

Furthermore, when Qwen3-VL raised an alarm its description of the fault
was correct on every single run, 90 of 90 Stage-3 \texttt{MATCH}, and it
never retracted a correct alarm at Stage~2. Its one weakness is a tendency
to fall silent with an \texttt{UNKNOWN} verdict on the hardest clip rather
than guess. Llama's behaviour is the concerning one, since on the log-less
clips its Stage~2 step does not merely mis-judge, it fabricates the
telemetry it was not given and then reasons from the fabrication, and a
monitor that invents a clean log to justify a \texttt{NORMAL} verdict is
worse than one that simply misses the fault.

Detecting an anomaly and describing it correctly are different abilities.
Qwen3-VL did both on every one of its positive runs, while GPT-4o did both
on about three quarters of its positive runs but tends to report a visible
hand instead of the underlying fault. Llama, although it fires on the
clips that include the log, produced a trustworthy account on a little
over half of its positive runs and collapses once the log is removed.

% ======================================================================
\section{Test 3: False-Alarm Behaviour}
\label{sec:test3}
% ======================================================================

\subsection{False-Alarm Test Results}

Unlike the first two tests, this one uses no anomaly clips at all. It
simulates the live monitor running against a single ten-minute recording
of normal robot operation. A window slides across the whole recording, one
analysis per window, and each analysis samples nine frames from inside its
window and runs the full two-stage pipeline. The window is fifteen seconds
long and advances fifteen seconds at a time, so the forty windows cover
the ten minutes without overlap and every second of the recording is seen
by exactly one analysis. Each of the three language models therefore
performs forty analyses over three hundred and sixty Stage-1 frames, which
is a realistic monitoring cadence of one analysis every fifteen seconds,
or two hundred and forty analyses per hour. On this material a verdict of
\texttt{NORMAL} is the correct one and every \texttt{ANOMALY} is, in
principle, a false alarm.

An earlier version of this test, which sampled a few random batches of
frames rather than sweeping the recording, showed all three models
inventing a human hand whenever the robot moved the yellow blocks. For the
run reported here an anti-hallucination clause was prepended to the
prompt. It tells the model that the robot itself picks and places the
yellow blocks as part of normal operation, that blocks changing position
between frames is not evidence of a human hand, and that an anomaly should
be reported only if a hand, finger, arm, person, or a clearly foreign
object is actually visible in one of the frames. YOLO was not included,
because a person-detecting rule would fire on any operator who walked
through the background, which is not the behaviour under study here.

\subsection{Results}

The results are summarised in Table~\ref{tab:test3}. GPT-4o and Qwen3-VL
each raised exactly one alarm across their forty analyses, both on the
same window (window~14, covering 195 to 210 seconds). Llama raised no
alarm at all. No model produced an \texttt{UNKNOWN} verdict, which is a
change from the earlier random-batch run where Qwen3-VL had degenerated
into unparseable output twice. On the two alarms that were raised, the
Stage~2 dense re-check confirmed the verdict rather than retracting it.

\begin{table}[h]
\centering
\small
\setlength{\tabcolsep}{6pt}
\begin{tabular}{@{}l c c c c c@{}}
\toprule
Model & Analyses & Stage-1 frames & Alarms raised & \texttt{UNKNOWN}
      & Mean time / analysis \\
\midrule
GPT-4o      & 40 & 360 & 1 (window 14) & 0 & 3\,s \\
Qwen3-VL 8B & 40 & 360 & 1 (window 14) & 0 & 11\,s \\
Llama-V 11B & 40 & 360 & 0             & 0 & 36\,s \\
\bottomrule
\end{tabular}
\caption{Test~3, continuous monitoring of ten minutes of normal operation.
Forty non-overlapping fifteen-second windows per model, nine frames
sampled per window, anti-hallucination clause in the prompt.}
\label{tab:test3}
\end{table}
\FloatBarrier

\subsection{The one alarm is a real hand}

The single window that GPT-4o and Qwen3-VL both flagged is not a
hallucination. A frame-by-frame look at window~14 shows a human hand and
forearm reaching into the workspace from the lower left between roughly
197 and 202 seconds, picking up and moving one of the yellow blocks, and
then withdrawing. The recording labelled as normal operation is therefore
not completely clean. It contains one brief human intrusion, and window~14
is the only window that overlaps it.

All three models saw the hand. GPT-4o's Stage-1 description reads ``a
human hand in the work area'' in two of the nine frames, and its verdict
is \texttt{ANOMALY}. Qwen3-VL's description reads ``a human hand in the
lower right corner'' and its verdict is also \texttt{ANOMALY}. Llama's
description reads ``a hand in Frame 5'', but its verdict is
\texttt{NORMAL}, on the reasoning that the hand ``is not in contact with
the robot or the blocks''. GPT-4o and Qwen3-VL were therefore correct on
window~14. Llama detected the same hand and then argued itself out of
reporting it, which on this window is a missed detection rather than a
clean pass.

\subsection{Conclusion for Test 3}

Once window~14 is set aside as a genuine intrusion, there are thirty-nine
windows of true normal operation. On those thirty-nine windows GPT-4o,
Qwen3-VL and Llama each raised zero alarms. The effective false-positive
rate over ten minutes of clean footage is therefore zero for all three
models, at a monitoring cadence of one analysis every fifteen seconds.
The anti-hallucination clause is what makes this possible. Without it, the
earlier random-batch run produced a hallucinated hand on roughly one
analysis in six. With it, the models correctly attribute block movement to
the robot and reserve an anomaly verdict for a visible intruder or object.

The one window that did contain a hand separates the models in a way that
matches the anomaly tests. GPT-4o and Qwen3-VL reported it, and their
Stage~2 re-check confirmed it. Llama saw the hand, described it, and then
suppressed the alarm with a criterion the task does not use, that a hand
must be touching the robot to count. For a safety monitor this is the
worse error, since it turns a real intrusion into a \texttt{NORMAL}
verdict.

Two limitations remain. The clean sample is a single ten-minute recording,
so the zero false-positive figure should be read as indicative rather than
as a measured rate, and a larger set of confirmed-normal footage would be
needed to put a real number on it. The anti-hallucination clause is
specific to this cell, since it names the yellow blocks and the robot's
pick-and-place operation, and it would have to be rewritten for a
different workspace.

% ======================================================================
\section{Discussion and general conclusions}
\label{sec:discussion}
% ======================================================================

\subsection{Ranking the detectors}

\textbf{GPT-4o} is the strongest and most stable detector overall. It
reached full Stage-1 recall on the common-anomaly set and 96\,\% on the
operational set, produced the most trustworthy justifications on the
external track, with three quarters of its positives graded MATCH, and was
by a wide margin the fastest of the language models, at a median of about
six seconds per run. Its weaknesses are specific and consistent. The
Stage~2 re-check repeatedly talks it out of correct detections of static
objects, and it under-describes composite anomalies, reporting the human
hand in a mixed clip but not the workpiece removal that accompanies it.

\textbf{Qwen3-VL 8B} is competitive on raw detection, level with GPT-4o at
Stage~1 on both tracks apart from the bottle clip, and it is the best
model at describing an operational fault. Every one of its
operational-anomaly positives was graded MATCH, including those made from
images alone, and it never retracted a correct operational alarm at
Stage~2. Against this, it is a thinking model that has been coerced into a
format it does not naturally produce, and it shows. Twelve of its runs
across the battery ended in \texttt{UNKNOWN}, and one generation failure,
the unfilled template, silently erased an entire clip's worth of
detections. It is also slow, at a median of about thirty-six seconds.

\textbf{Llama 3.2 Vision 11B} is the weakest of the three. It is blind to
the two external object-only clips, it loses more than three quarters of
its operational-anomaly recall when the state log is removed, and even
when it does fire its account is usually a hallucinated scene with a hand
mentioned somewhere inside it, so fewer than one external positive in five
was graded MATCH. Its Stage~2 behaviour on log-less clips, in which it
fabricates a clean log to justify \texttt{NORMAL}, is the most dangerous
single behaviour observed. It is also the slowest model, at a median of
about eighty seconds, because the caption-then-synthesise pipeline makes
one model call per frame.

\textbf{YOLO11n-pose} matched the language models' recall on human
intrusions, at 90\,\% on the external set, at a thousandth of the latency,
and was the only detector to recover the one-frame hand in an1. But it is
structurally incapable of the rest of the task. It cannot see a foreign
object, it fired spuriously on the bottle clip, and it has no access to
and no concept of the operational channel. Counting both tracks, it
identified the correct cause in only about half of its detections. The
detection-only variant (\texttt{yolo\_coco}) is weaker still, at 70\,\%
recall on the external set.

\subsection{The role of the state log}

The clearest quantitative result in the battery is the effect of the
machine-state log on operational-anomaly detection. With the log, all
three language models detected every fault in every run. Without it,
GPT-4o and Qwen3-VL held about 80\,\% while Llama fell to 22\,\%. The log
is a very strong prior, since it names the fault explicitly, and the
with-log condition should be read as a control rather than as a result.
The finding of substance is that the two stronger models degrade
gracefully rather than catastrophically when the log is withdrawn, which
suggests that a vision--language monitor could provide useful
operational-anomaly coverage as a fallback when controller telemetry is
missing or untrusted. The exception is the fault in which the robot skips
its pick altogether, which produces essentially no visual evidence and
was missed by every model without the log.

\subsection{The two-stage pipeline}

As configured for this battery, the Stage~2 re-check is a net liability on
the anomaly sets. It retracted 39 correct detections across the two sets,
seventeen for GPT-4o, ten for Qwen3-VL and twelve for Llama, and made
exactly one useful correction, on the false-alarm video. This is an
artefact of the evaluation design as much as of the pipeline, because with
no normal clips in the anomaly sets a retraction there is always wrong, so
any overturning mechanism can only lose. The retractions themselves fall
into three groups. The first is genuine second-thoughts, in which the
model looks again at a real object or a real empty gripper and reasons its
way to \texttt{NORMAL}, and these are roughly half. The second is
mechanical parser failures, in which a broken or templated model output is
read as a \texttt{NORMAL} verdict. The third is localisation failures, in
which the dense window is placed over a part of the clip that does not
contain the transient event.

The recommendation that follows is to keep Stage~2 for what it does well,
confirming an event against denser evidence and localising it in time, but
to disable its authority to overturn the Stage~1 verdict until the
evaluation includes enough normal material to show that overturning is, on
balance, beneficial. The localisation output itself remained useful. Where
the anchor was placed correctly, the reported anomaly window overlapped
the true window in essentially every case.

\subsection{The character of the failures}

A recurring theme across the three tests is that the language models fail
by incompleteness and by hallucination, but not by confusion. Across 600
runs on clips that genuinely contained an anomaly, Stage~3 did not grade a
single justification MISMATCH. No model ever looked at a clip with a hand
in it and reported a bottle, or looked at a bottle and reported a hand.
When these models are wrong on a real anomaly, they have either described
it too vaguely, which is PARTIAL, or not fired at all. The one place a
fabricated anomaly appeared was the earlier random-batch version of the
false-alarm test, where all three models invented a human hand to explain
the robot moving its own workpieces. The anti-hallucination clause added
for the continuous version removed that failure, and in the continuous run
the only alarm raised on the normal recording turned out to be a real hand
that the recording was not supposed to contain. For a monitoring
application this is a reassuring shape of error on the detection side, and
on the false-alarm side it shows the behaviour is controllable through the
prompt.

\subsection{Limitations}

The anomaly sets contain no negative controls, so all figures on them are
recall and no precision or F-measure can be computed. The false-alarm
test runs on a single ten-minute recording, so its zero false-positive
rate over thirty-nine clean windows should be treated as indicative and
not as a measured rate. The operational no-log clips are re-used footage from the
with-log clips, which isolates the log's contribution cleanly but means
the no-log results are not from independent recordings. Qwen3-VL was run
in a forced non-thinking mode that it does not natively support, and some
of its \texttt{UNKNOWN} outcomes may be an artefact of that coercion
rather than of the vision task. Finally, the Stage~3 judge is itself
GPT-4o. Its grades of GPT-4o's own justifications should be read with that
in mind, although the judge is given only the two short texts and not the
identity of the detector.

% ======================================================================
\appendix
\section{Latency}

\begin{table}[h]
\centering
\begin{tabular}{@{}l l c c c c@{}}
\toprule
Model & Track & Mean (s) & Median (s) & Min (s) & Max (s) \\
\midrule
GPT-4o          & external & 8.8  & 6.3  & 5.3  & 26.3 \\
GPT-4o          & internal & 10.4 & 6.7  & 3.2  & 48.7 \\
Qwen3-VL 8B     & external & 36.2 & 35.7 & 22.8 & 46.1 \\
Qwen3-VL 8B     & internal & 41.1 & 44.6 & 22.9 & 46.8 \\
Llama-V 11B     & external & 71.1 & 80.0 & 34.2 & 82.7 \\
Llama-V 11B     & internal & 68.2 & 81.4 & 34.7 & 84.4 \\
YOLO11n (either)& external & $\approx 2$ & $\approx 1.5$ & 0.1 & 4.6 \\
\bottomrule
\end{tabular}
\caption{Wall-clock time per run (Stage~1 and Stage~2 combined). GPT-4o
over the API is four to twelve times faster than the local language
models, and YOLO is two to three orders of magnitude faster than any of
them. The complete battery took approximately GPT-4o 34~min, Qwen3-VL
131~min, Llama 234~min, YOLO under a minute, and the false-alarm test
9~min.}
\label{tab:latency}
\end{table}

\section{Data and reproduction}

All results are derived from four comma-separated files and the per-run
transcripts. The file \texttt{summary.csv} holds the 600 language-model
runs, with per-run Stage-1 and Stage-2 verdicts, the final verdict, the
localised window and the latency. The file \texttt{yolo\_summary.csv}
holds the 200 YOLO rows. The file \texttt{false\_alarm.csv} holds the 18
false-alarm runs, each recording the exact frame indices that were
sampled. The file \texttt{stage3.csv} holds the 618 semantic judgements.
The full Stage-1 and Stage-2 text of every run is kept under
\texttt{eval\_results\_experiment/}. The reference anomaly for each clip
is defined in \texttt{ground\_truth.py}. The whole battery is reproduced
by \texttt{run\_battery.sh}, and the tables in this report are regenerated
by \texttt{report\_battery.py}, \texttt{score\_vs\_truth.py} and
\texttt{stage3\_match.py}.

\end{document}
